\documentclass[10pt,a4paper]{amsart}
\usepackage[margin=19mm]{geometry}
\usepackage{amsmath,amssymb,mathtools}
\usepackage[scr=rsfs,cal=euler]{mathalfa}
\usepackage{xfrac}
\usepackage[unicode,hidelinks,bookmarks=false]{hyperref}
\newcommand{\ie}{i.e.\ }
\newcommand{\cf}{cf.\ }
\newcommand{\resp}{respectively\ }
\newcommand{\Subsection}{\subsection}
\providecommand{\nobreakdash}{\nobreak\hbox{-}}
\setlength{\emergencystretch}{2em}
\multlinegap0pt
\usepackage[shortlabels]{enumitem}
\usepackage{amssymb}
\usepackage{tikz}
\newtheorem{lem}{Lemma}
\newtheorem{defn}{Definition}
\newtheorem{fact}{Fact}
\newtheorem{claim}{Claim}
\newcommand{\EE}{\mathbb{E}}
\newcommand{\RR}{\mathbb{R}}
\DeclareMathOperator{\Var}{Var}
\newcommand{\paren}[1]{\left( #1 \right)}
\title{On the largest sum-free subset of the lattice cube}
\author{Peter Keevash, Jeck Lim}
\date{Source: arXiv:2605.00816v2 (CC BY 4.0)}
\begin{document}
\maketitle

\noindent\emph{Source and licence.} On the largest sum-free subset of the lattice cube. Version arXiv:2605.00816v2 (CC BY 4.0), distributed by arXiv under the Creative Commons Attribution 4.0 International licence, http://creativecommons.org/licenses/by/4.0/, as recorded in the arXiv licence metadata for this exact version and shown on the abstract page https://arxiv.org/abs/2605.00816v2. Full licence notice: \texttt{licenses/arxiv-2605.00816-CC-BY-4.0.txt}.\par\smallskip

\noindent\textit{Editorial scope.} Counted unit: the article's lem lem:real-layer-estimate (source0.tex lines 287--293). The article's complete original proof of it is delivered with it (source0.tex lines 359--441, one \texttt{proof} environment of 83 lines, which the article places away from the statement). No mathematics is written, repaired or re-derived: the statement and the proof are the article's own lines, in the article's own notation, and no retained line is substituted or re-flowed.  Retained uncounted as the source-local context the proof needs: lines 322--327; lines 331--334; lines 656--656.  Nothing was omitted from any retained span, so the package carries no omission record.  No retained line contains a citation, so the wrapper carries no bibliography and no bibliography entry.  No retained line contains a figure environment, an image-inclusion command or drawing code, so no image is delivered and the package declares no asset dependency.  Editorial formatting only, in addition: an AMS \texttt{amsart} preamble that restates the article's own declarations and macros used by the retained lines, namely the environments \texttt{lem}, \texttt{defn}, \texttt{fact}, \texttt{claim} on separate counters, together with the standard packages those lines need.  The retained environments print in this document as Lemma 1, Definition 1, Fact 1, Claim 1 in order of appearance, whereas the article prints the same items with the numbering of its own full document.  The complete native source of the article as distributed in the arXiv e-print for this exact version is bundled unchanged as \texttt{sources/199596/source0.tex}.
\begin{defn}
    A distribution on $\RR$ with probability density function $f:\RR\to [0,\infty)$ is \emph{log-concave} if $\log f$ is a concave function.
\end{defn}
\begin{fact} \label{fact:log-concave-add}
    If $X,Y$ are independent log-concave random variables, then $X+Y$ is also log-concave.
\end{fact}

\begin{lem} \label{lem:unimodal}
    Let $X$ be a random variable on $\RR$ with probability density function $f$. Suppose $\EE[X]=\mu, \Var(X)=\sigma^2$, and $a\in \RR$ is such that $f$ is increasing on $(-\infty,a)$ and decreasing on $(a,\infty)$. Then 
    \[|a-\mu|\leq \sqrt{3}\sigma.\]
\end{lem}

\section{Numerical verification for small \texorpdfstring{$d$}{d}} \label{sec:numerical}

\begin{lem} \label{lem:real-layer-estimate}
    The following inequalities hold for $d\geq 3$.
    \begin{enumerate}
        \item $f_{d-1}(2u_d^*-1)\geq 2.01f_{d-1}(u_d^*-1).$
        \item $0.99f_{d-1}(u_d^*)\geq 2f_{d-1}(u_d^*-1)+4f_{d-1}(2u_d^*).$
    \end{enumerate}
\end{lem}

\begin{proof}[Proof of Lemma~\ref{lem:real-layer-estimate}]
    Let $X$ be the uniform random variable on $[0,1]$, then $f_d$ is the probability density function of the sum of $d$ independent copies of $X$. Since the probability density function of $X$ is the indicator function on $[0,1]$, which is log-concave, $f_d$ is also log-concave by Fact~\ref{fact:log-concave-add}. In particular, $f_d'(x)/f_d(x)$ is a decreasing function of $x$. Let $M(t)=\EE[e^{tX}]=\frac{e^t-1}{t}$ be the moment generating function of $X$.

    For $\sigma\in\RR$, let $X_\sigma$ be the exponential tilt of $X$, which is the random variable with probability density function
    \[f(x;\sigma) = \frac{e^{\sigma x}}{M(\sigma)}f(x).\]
    This is indeed a probability density function since
    \[\int_{-\infty}^\infty f(x;\sigma)dx = \frac{1}{M(\sigma)}\int_{-\infty}^{\infty} e^{\sigma x}f(x)dx = 1.\]
    Let us compute the mean and variance of $X_\sigma$. We have
    \begin{align*}
        \EE[X_\sigma] &= \int_{-\infty}^\infty xf(x;\sigma)dx = \frac{1}{M(\sigma)}\int_{-\infty}^\infty xe^{\sigma x}f(x)dx = \frac{M'(\sigma)}{M(\sigma)}, \\
        \EE[X_\sigma^2] &= \int_{-\infty}^\infty x^2f(x;\sigma)dx = \frac{1}{M(\sigma)}\int_{-\infty}^\infty x^2e^{\sigma x}f(x)dx = \frac{M''(\sigma)}{M(\sigma)}.
    \end{align*}
    Therefore, 
    \[\Var(X_\sigma) = \EE[X_\sigma^2] - \EE[X_\sigma]^2 = \frac{M''(\sigma)}{M(\sigma)} - \paren{\frac{M'(\sigma)}{M(\sigma)}}^2.\]

    Let $Y_\sigma$ be the sum of $d$ independent copies of $X_\sigma$, and let $f_d(\cdot;\sigma)$ be the probability density function of $Y_\sigma$. Since $f(\cdot;\sigma)$ is log-concave, $f_d(\cdot;\sigma)$ is also log-concave by Fact~\ref{fact:log-concave-add}. 

    \begin{claim}
        We have
        \[f_d(x;\sigma) = \frac{e^{\sigma x}}{M(\sigma)^d}f_d(x).\]
    \end{claim}
    \begin{proof}
        We argue by induction on $d$. For $d=1$, this is true by definition. For $d\geq 2$, we have
        \begin{align*}
            f_d(x;\sigma) &= \int_{-\infty}^\infty f_{d-1}(x-y;\sigma)f(y;\sigma)dy \\
            &= \int_{-\infty}^\infty \frac{e^{\sigma (x-y)}}{M(\sigma)^{d-1}}f_{d-1}(x-y) \cdot \frac{e^{\sigma y}}{M(\sigma)}f(y) dy \\
            &= \frac{e^{\sigma x}}{M(\sigma)^d} \int_{-\infty}^\infty f_{d-1}(x-y)f(y) dy = \frac{e^{\sigma x}}{M(\sigma)^d}f_d(x).
        \end{align*}
    \end{proof}

    We would like to show that $f_d'(x)/f_d(x)$ is approximately $-\sigma^*$ for $x$ close to $d/3$, where $\sigma^*=-2.149...$ is the unique solution to $M'(\sigma^*)/M(\sigma^*)=1/3$. By log-concavity, $f_d'/f_d$ is decreasing in $x$. We will show a lower bound on $f_d'(x)/f_d(x)$ for $x$ slightly above $d/3$, and an upper bound for $x$ slightly below $d/3$. For our application, we are aiming for the lower and upper bounds of $\sigma_1:=-3$ and $\sigma_2:=-1.6$ respectively.

    Let $\mu_i=\EE[X_{\sigma_i}]$ and $s_i^2=\Var(X_{\sigma_i})$ for $i=1,2$. We have the following approximate values:
    \begin{align*}
        \mu_1 &\approx 0.280938, \quad s_1 \approx 0.236580, \\
        \mu_2 &\approx 0.372030, \quad s_2 \approx 0.271405.
    \end{align*}
    We have $\EE[Y_{\sigma_i}] = d\mu_i$ and $\Var(Y_{\sigma_i}) = d s_i^2$. Since $f_d(\cdot;\sigma_i)$ is log-concave, it is unimodal, so by Lemma~\ref{lem:unimodal}, we have $f_d'(x;\sigma_i)\leq 0$ for $x\geq d\mu_i + \sqrt{3d} s_i$, and $f_d'(x;\sigma_i)\geq 0$ for $x\leq d\mu_i - \sqrt{3d} s_i$. Therefore, for $x\geq d\mu_i + \sqrt{3d} s_i$, we have
    \begin{align*}
        0 &\geq f_d'(x;\sigma_i) = \frac{e^{\sigma_i x}}{M(\sigma_i)^d}(\sigma_i f_d(x) + f_d'(x)) \\
        \implies \frac{f_d'(x)}{f_d(x)} &\leq -\sigma_i.
    \end{align*}
    Similarly, $f_d'(x)/f_d(x)\geq -\sigma_i$ for $x\leq d\mu_i - \sqrt{3d} s_i$. Thus, we have
    \[-\sigma_2 \leq \frac{f_d'(x)}{f_d(x)} \leq -\sigma_1\]
    for $x\in [d\mu_1 + \sqrt{3d} s_1, d\mu_2 - \sqrt{3d} s_2]$. In particular, this implies that for $d\mu_1 + \sqrt{3d} s_1 \leq x\leq y\leq d\mu_2 - \sqrt{3d} s_2$, we have
    \begin{equation} \label{eq:fd-ratio-bound}
        e^{-\sigma_2(y-x)} \leq \frac{f_d(y)}{f_d(x)} \leq e^{-\sigma_1(y-x)}.
    \end{equation}
    For $d\geq 200$, the interval $[d\mu_1 + \sqrt{3d} s_1, d\mu_2 - \sqrt{3d} s_2]$ contains $[d/3-1, d/3+1]$, so \eqref{eq:fd-ratio-bound} holds for all $x,y\in [d/3-1, d/3+1]$. For $5\leq d < 200$, we directly verify the lemma by numerically computing the relevant quantities, see Appendix~\ref{sec:numerical}.

    We now estimate the value of $u_d^*$. Let $u_d^* = d/3 + \alpha$, then $f_d(u_d^*)=2f_d(2u_d^*)$ implies that
    \[\frac{f_d(d/3+\alpha)}{f_d(2d/3+2\alpha)} = \frac{f_d(d/3+\alpha)}{f_d(d/3-2\alpha)} = 2.\]
    \eqref{eq:fd-ratio-bound} then gives
    \begin{align*}
        e^{-3\sigma_2\alpha} &\leq 2\leq e^{-3\sigma_1\alpha}\\
        \implies 0.077016 &\leq \alpha \leq 0.144406.
    \end{align*}
    Note that a priori \eqref{eq:fd-ratio-bound} may not apply since we do not know the value of $\alpha$ yet. However, the above estimates show that $f_d(d/3+0.5)/f_d(d/3-1)>2$. Then, the uniqueness of $\alpha$ implies that $\alpha<0.5$, so \eqref{eq:fd-ratio-bound} does apply.

    Recall that the lemma is the following two inequalities.
    \begin{enumerate}
        \item $f_{d-1}(2u_d^*-1)\geq 2.01f_{d-1}(u_d^*-1).$
        \item $0.99f_{d-1}(u_d^*)\geq 2f_{d-1}(u_d^*-1)+4f_{d-1}(2u_d^*).$
    \end{enumerate}
    In our computations below, we only evaluate $f_{d-1}(x)$ for $x\in [(d-1)/3-1, (d-1)/3+1]$, so \eqref{eq:fd-ratio-bound} holds.

    For (1), we have
    \begin{align*}
        \frac{f_{d-1}(2u_d^*-1)}{f_{d-1}(u_d^*-1)} &= \frac{f_{d-1}(2d/3+2\alpha-1)}{f_{d-1}(d/3+\alpha-1)} \\
        &= \frac{f_{d-1}(d/3-2\alpha)}{f_{d-1}(d/3+\alpha-1)} \\
        &\geq e^{-\sigma_2(1-3\alpha)}\\
        &\geq e^{-\sigma_2}/2 \approx 4.953032/2 > 2.01.
    \end{align*}

    For (2), we have
    \begin{align*}
        \frac{2f_{d-1}(u_d^*-1)}{f_{d-1}(u_d^*)} + \frac{4f_{d-1}(2u_d^*)}{f_{d-1}(u_d^*)} &= \frac{2f_{d-1}(d/3+\alpha-1)}{f_{d-1}(d/3+\alpha)} + \frac{4f_{d-1}(2d/3+2\alpha)}{f_{d-1}(d/3+\alpha)} \\
        &= \frac{2f_{d-1}(d/3+\alpha-1)}{f_{d-1}(d/3+\alpha)} + \frac{4f_{d-1}(d/3-2\alpha-1)}{f_{d-1}(d/3+\alpha)} \\
        &\leq 2e^{\sigma_2}+4e^{\sigma_2(3\alpha+1)} \\
        &\leq 0.961800 < 0.99.
    \end{align*}

\end{proof}

\end{document}
